\documentclass[../DoAn.tex]{subfiles}
\begin{document}

% ═══════════════════════════════════════════════════════════════════════════════
\section{Thiết kế Kiến trúc Tổng thể}
\label{sec:architecture}
% ═══════════════════════════════════════════════════════════════════════════════

Hệ thống được xây dựng theo mô hình ứng dụng phi tập trung (Decentralized Application -- DApp) với kiến trúc lai ba tầng giữa thành phần on-chain, off-chain và client, nhằm cân bằng giữa tính phi tập trung, hiệu năng xử lý và trải nghiệm người dùng~\cite{erc4337,groth16}.

\begin{figure}[H]
    \centering
    \includegraphics[width=1\textwidth]{Hinhve/overall_architecture.png}
    \caption{Kiến trúc tổng thể ba tầng của hệ thống Stealth Wallet}
    \label{fig:ArchitectureDiagram}
\end{figure}

\subsection{Tầng Smart Contract (On-chain Layer)}
\label{subsec:onchain_layer}

Tầng này đảm nhiệm việc lưu trữ trạng thái không thể thay đổi và xác minh các bằng chứng mật mã trên blockchain Ethereum Sepolia. Các hợp đồng thông minh được triển khai bao gồm:

\begin{itemize}
    \item \textbf{\texttt{StealthTreeManager}:} Lưu trữ \texttt{Merkle Root} hiện hành của cây Merkle Poseidon độ sâu 20. Hợp đồng tích hợp trực tiếp với \texttt{SMTUpdateVerifier} để xác minh bằng chứng Groth16 trước mỗi lần cập nhật gốc cây. Đây là điểm đồng thuận duy nhất giữa trạng thái off-chain (server) và on-chain.

    \item \textbf{\texttt{StealthAccountFactory}:} Factory triển khai các \texttt{StealthAccount} thông qua cơ chế \texttt{CREATE2}. Địa chỉ stealth được xác định trước dựa trên \texttt{indexCommitment} đóng vai trò là \texttt{salt}, cho phép tài sản gửi đến tài khoản trước khi hợp đồng được triển khai thực sự~\cite{eip1014}.

    \item \textbf{\texttt{StealthAccount}:} Smart Account tương thích ERC-4337, tích hợp:
    \begin{itemize}
        \item Hàm \texttt{validateUserOp()}: xác thực \texttt{UserOperation} thông qua bằng chứng Groth16 ZK.
        \item Hàm \texttt{execute()}: thực thi lệnh chuyển tài sản thông thường sau khi xác minh ZK.
        \item Các hàm \texttt{executeStealthTransfer()}, \texttt{executeERC20StealthTransfer()}, \texttt{executeERC721StealthTransfer()}: thực thi chuyển khoản kèm phát sự kiện \texttt{Announcement}.
    \end{itemize}

    \item \textbf{\texttt{SocialRecoveryModule}:} Hợp đồng độc lập per-account, quản lý vòng đời phục hồi với ba hàm chính: \texttt{proposeRecovery()}, \texttt{approveRecovery()}, \texttt{executeRecovery()}, liên kết trực tiếp với chỉ mục lá trong Merkle Tree thông qua \texttt{mappedIndex}.

    \item \textbf{\texttt{ERC5564Announcer}:} Phát sự kiện \texttt{Announcement(schemeId, stealthAddress, caller, ephemeralPubKey, metadata)} theo chuẩn ERC-5564~\cite{erc5564}, phục vụ quá trình quét địa chỉ stealth của người nhận.

    \item \textbf{\texttt{StealthPaymaster}:} Hợp đồng Paymaster theo ERC-4337, tài trợ phí Gas cho mọi \texttt{UserOperation} nhằm giải quyết bài toán Cold Start của Stealth Account mới tạo.

    \item \textbf{\texttt{StealthSpendVerifier}} và \textbf{\texttt{SMTUpdateVerifier}:} Hai hợp đồng xác minh bằng chứng Groth16 được xuất tự động từ mạch Circom thông qua SnarkJS. \texttt{StealthSpendVerifier} dùng cho mạch chi tiêu (spend), \texttt{SMTUpdateVerifier} dùng cho mạch cập nhật cây.
\end{itemize}

\subsection{Tầng Relayer Server (Off-chain Layer)}
\label{subsec:server_layer}

Máy chủ Node.js đóng vai trò là bộ điều phối trung tâm, chịu trách nhiệm duy trì trạng thái cây và tổng hợp \texttt{UserOperation}. Tầng này hoạt động hoàn toàn ngoài chuỗi nhằm giảm chi phí tính toán Gas và tăng tốc độ phản hồi:

\begin{enumerate}
    \item \textbf{Quản lý Merkle Tree:} Server duy trì một bản sao \texttt{PoseidonSparseMerkleTree} hoàn chỉnh cục bộ, được khởi tạo từ tệp \texttt{leaves.json} và đồng bộ với trạng thái on-chain khi khởi động. Cập nhật cây sử dụng thuật toán sparse với độ phức tạp $O(d)$, chỉ tính lại các nút trên đường từ lá đến gốc.

    \item \textbf{Sinh bằng chứng ZK:} Server gọi thư viện SnarkJS để sinh bằng chứng Groth16 cho cả hai mạch: mạch chi tiêu (\texttt{stealth\_spend}) và mạch cập nhật cây (\texttt{smt\_update}). Các tệp \texttt{.wasm} và \texttt{.zkey} được nạp sẵn trong bộ nhớ để giảm trễ.

    \item \textbf{Tổng hợp UserOperation:} Sau khi có bằng chứng ZK, server xây dựng \texttt{UserOperation} hoàn chỉnh và gọi \texttt{entryPoint.handleOps()} trực tiếp thông qua ví quản lý (signer) của relayer.

    \item \textbf{Cơ chế Rollback:} Nếu giao dịch on-chain thất bại sau khi lá đã được chèn vào cây cục bộ, hàm \texttt{rollbackLastLeaf()} được gọi để hoàn tác thay đổi và duy trì tính nhất quán.
\end{enumerate}

\subsection{Tầng Client (Frontend Layer)}
\label{subsec:client_layer}

Giao diện React.js + Ethers.js v6 đảm nhiệm tất cả phép tính mật mã phía người dùng, đảm bảo khóa bí mật không bao giờ rời khỏi trình duyệt của người dùng:

\begin{itemize}
    \item Sinh và lưu trữ cặp khóa (scanPriv, spendPriv) cục bộ trong bộ nhớ phiên.
    \item Tính toán ECDH shared secret và dẫn xuất địa chỉ stealth.
    \item Quét sự kiện \texttt{Announcement} từ blockchain để phát hiện tài sản đến.
    \item Ủy quyền sinh bằng chứng ZK và gửi \texttt{UserOperation} cho server Relayer.
\end{itemize}

% ═══════════════════════════════════════════════════════════════════════════════
\section{Thiết kế Chi tiết Các Thành phần}
\label{sec:detailed_design}
% ═══════════════════════════════════════════════════════════════════════════════

\subsection{Thiết kế luồng Tạo Ví và Đăng ký Lá}
\label{subsec:create_wallet_flow}

Khi người dùng khởi tạo ví stealth mới, hệ thống thực hiện tuần tự các bước sau:

\begin{enumerate}
    \item \textbf{[Client]} Sinh cặp khóa ngẫu nhiên:
    \[
        (\texttt{scanPriv}, v) \xleftarrow{\$} \mathbb{Z}_n^*, \quad V = v \cdot G
    \]
    \[
        (\texttt{spendPriv}, x) \xleftarrow{\$} \mathbb{F}_p
    \]

    \item \textbf{[Client]} Tính giá trị lá Merkle:
    \[
        k = \text{Poseidon}(x)
    \]
    và địa chỉ danh tính:
    \[
        \texttt{addr} = \text{computeAddress}(\texttt{spendPub})
    \]

    \item \textbf{[Client $\to$ Server]} Gửi yêu cầu \texttt{POST /leaves} với payload \(\{k, \texttt{addr}\}\).

    \item \textbf{[Server]} Kiểm tra đồng bộ (\texttt{isTreeSynced}), kiểm tra trùng lá (\texttt{hasLeaf(k)}):
    \begin{itemize}
        \item Nếu lá đã tồn tại: trả về index và indexCommitment hiện có.
        \item Nếu chưa: gán chỉ mục \(i = \texttt{nextIndex}\), lấy Merkle proof của slot trống, chèn lá vào cây.
    \end{itemize}

    \item \textbf{[Server $\to$ Blockchain]} Sinh bằng chứng Groth16 SMT Update và gọi:
    \[
        \texttt{contract.updateRoot}(\texttt{newRoot}, k, i, \texttt{auth})
    \]
    Nếu giao dịch on-chain thất bại, \texttt{rollbackLastLeaf()} được gọi.

    \item \textbf{[Server $\to$ Client]} Trả về \(\{i, \texttt{newRoot}, \texttt{txHash}\}\).

    \item \textbf{[Client]} Tính:
    \[
        \texttt{IndexHash} = \text{Poseidon}(i, 0)
    \]
    và lưu bộ meta-address đầy đủ: \((\texttt{scanPriv}, \texttt{spendPriv}, V, \texttt{spendPub}, i, \texttt{IndexHash})\).
\end{enumerate}

\begin{figure}[H]
    \centering
    \includegraphics[width=0.8\textwidth]{Hinhve/usecase_create_wallet.png}
    \caption{Biểu đồ Use Case -- Chức năng Tạo ví và Đăng ký lá}
    \label{fig:uc_create_wallet}
\end{figure}

\begin{figure}[H]
    \centering
    \includegraphics[width=1\textwidth]{Hinhve/class_create_wallet.png}
    \caption{Biểu đồ Lớp -- Chức năng Tạo ví và Đăng ký lá}
    \label{fig:class_create_wallet}
\end{figure}

\begin{figure}[H]
    \centering
    \includegraphics[width=1\textwidth]{Hinhve/create_wallet_sq_diagram.png}
    \caption{Biểu đồ tuần tự UML -- Luồng tạo ví và đăng ký lá Merkle}
    \label{fig:sd_create_wallet}
\end{figure}

\subsection{Thiết kế luồng Gửi Tài sản Stealth}
\label{subsec:send_flow}

Luồng gửi tài sản hoàn chỉnh từ một Stealth Account đã tồn tại bao gồm 9 bước, được tổng hợp trong Bảng~\ref{table:send_flow} và minh họa trong Hình~\ref{fig:sd_send_stealth}:

\begin{table}[H]
\centering
\resizebox{\textwidth}{!}{%
\begin{tabular}{|c|l|l|p{7cm}|}
\hline
\textbf{Bước} & \textbf{Bên thực hiện} & \textbf{Hành động} & \textbf{Chi tiết} \\ \hline
1 & Client & Tra cứu Meta-address & Đọc \texttt{scanPub}, \texttt{spendPub}, \texttt{IndexHash} của người nhận (từ ENS hoặc nhập tay) \\ \hline
2 & Client & Tính ECDH Shared Secret & $S = r \cdot V_{\text{Bob}}$; $h = \text{keccak256}(S.x \| S.y)$ \\ \hline
3 & Client & Tính indexCommitment & $C = \text{Poseidon}(\text{IndexHash} + h)$ \\ \hline
4 & Client & Tính địa chỉ stealth & $A_{\text{stealth}} = \text{Factory.getAddress}(C)$ \\ \hline
5 & Client $\to$ Server & Gửi yêu cầu chi tiêu & \texttt{POST /api/spend-zk-proof} với \texttt{spendPriv}, \texttt{sharedSecretHash}, \texttt{senderIndexCommitment} \\ \hline
6 & Server & Sinh bằng chứng ZK & \texttt{generateSpendProof()} gọi SnarkJS với mạch \texttt{stealth\_spend} \\ \hline
7 & Server & Xây dựng UserOperation & \texttt{callData} = \texttt{executeStealthTransfer()}, \texttt{signature} = ABI-encode(auth, indexCommitment) \\ \hline
8 & Server & Debug xác minh on-chain & Gọi \texttt{debugVerifyProof()} trên hợp đồng trước khi gửi \\ \hline
9 & Server $\to$ Blockchain & Gửi handleOps & \texttt{entryPoint.handleOps([userOp], beneficiary)} \\ \hline
\end{tabular}
}
\caption{Luồng xử lý giao dịch stealth chi tiết}
\label{table:send_flow}
\end{table}

\begin{figure}[H]
    \centering
    \includegraphics[width=0.8\textwidth]{Hinhve/send_stealth_payment}
    \caption{Biểu đồ Use Case -- Chức năng Gửi tài sản Stealth}
    \label{fig:uc_send_stealth}
\end{figure}

\begin{figure}[H]
    \centering
    \includegraphics[width=1\textwidth]{Hinhve/class_send_stealth_payment.png}
    \caption{Biểu đồ Lớp -- Chức năng Gửi tài sản Stealth}
    \label{fig:class_send_stealth}
\end{figure}

\begin{figure}[H]
    \centering
    \includegraphics[width=1\textwidth]{Hinhve/send_stealth_payment.png}
    \caption{Biểu đồ tuần tự UML -- Luồng gửi tài sản stealth (ETH/ERC-20/ERC-721)}
    \label{fig:sd_send_stealth}
\end{figure}

\subsection{Thiết kế luồng Quét Tài sản Stealth}
\label{subsec:scan_flow}

Sau khi nhận tài sản qua giao thức Stealth Address, người nhận (Bob) cần quét blockchain để phát hiện các tài khoản thuộc sở hữu của mình. Quy trình quét được mô tả trong Hình~\ref{fig:sd_scan_stealth}:

\begin{enumerate}
    \item \textbf{[Client]} Truy vấn toàn bộ sự kiện \texttt{Announcement} từ hợp đồng \texttt{ERC5564Announcer} trong phạm vi block được chỉ định.

    \item \textbf{[Client]} Với mỗi sự kiện, tái tạo Shared Secret bằng khóa xem:
    \[
        S' = v_{\text{Bob}} \cdot R, \quad h' = \text{keccak256}(S'.x \,\|\, S'.y)
    \]

    \item \textbf{[Client]} Tính \texttt{indexCommitment} và địa chỉ stealth kỳ vọng:
    \[
        C' = \text{Poseidon}(\text{IndexHash}_{\text{Bob}} + h'), \quad A_{\text{expected}} = \text{Factory.getAddress}(C')
    \]

    \item \textbf{[Client]} So sánh \(A_{\text{expected}}\) với \texttt{stealthAddress} trong sự kiện. Nếu khớp: ghi nhận tài khoản và lưu \(h'\) phục vụ sinh ZK proof sau này.
\end{enumerate}

\begin{figure}[H]
    \centering
    \includegraphics[width=0.8\textwidth]{Hinhve/usecase_scan_stealth_payment.png}
    \caption{Biểu đồ Use Case -- Chức năng Quét và phát hiện tài sản stealth}
    \label{fig:uc_scan_stealth}
\end{figure}

\begin{figure}[H]
    \centering
    \includegraphics[width=1\textwidth]{Hinhve/class_scan_wallet.png}
    \caption{Biểu đồ Lớp -- Chức năng Quét và phát hiện tài sản stealth}
    \label{fig:class_scan_stealth}
\end{figure}

\begin{figure}[H]
    \centering
    \includegraphics[width=0.95\textwidth]{Hinhve/scan_stealth_payment.png}
    \caption{Biểu đồ tuần tự UML -- Luồng quét và phát hiện tài sản stealth}
    \label{fig:sd_scan_stealth}
\end{figure}

\subsection{Thiết kế Cơ sở Dữ liệu Off-chain}
\label{subsec:offchain_db}

Server sử dụng tệp JSON (\texttt{leaves.json}) làm cơ sở dữ liệu ngoài chuỗi với cấu trúc như sau:

\begin{verbatim}
{
  "nextIndex": <number>,
  "leafMap": {
    "<index>": "0x<commitment_hex>",
    ...
  },
  "tree": {
    "depth": 20,
    "nodes": { ... }   // snapshot các nút non-zero
  }
}
\end{verbatim}

Mỗi trường có vai trò cụ thể:
\begin{itemize}
    \item \textbf{\texttt{nextIndex}:} Bộ đếm tăng đơn điệu, không bao giờ giảm kể cả sau khi xóa lá. Đảm bảo mỗi vị trí chỉ được gán một lần.
    \item \textbf{\texttt{leafMap}:} Ánh xạ thưa (sparse) từ chỉ mục sang giá trị cam kết lá $k$. Chỉ lưu các lá đang hoạt động (không lưu slot đã xóa).
    \item \textbf{\texttt{tree.nodes}:} Snapshot các nút nội bộ non-zero của cây Poseidon. Cho phép khởi động nhanh mà không cần tái tính toán toàn bộ cây từ các lá.
\end{itemize}

Thiết kế này giúp thời gian khởi động lại server từ $O(N \cdot d)$ (tái xây dựng đầy đủ) giảm xuống $O(1)$ (nạp snapshot).

\subsection{Thiết kế luồng Phục hồi Xã hội}
\label{subsec:recovery_flow}

Luồng phục hồi Social Recovery diễn ra qua bốn giai đoạn chính (Hình~\ref{fig:sd_social_recovery}):

\begin{enumerate}
    \item \textbf{Giai đoạn 1 -- Thiết lập Guardian (Chủ sở hữu):}  Triển khai hợp đồng \texttt{SocialRecoveryModule} với danh sách guardian và ngưỡng phê duyệt $t$. Liên kết \texttt{mappedIndex} = chỉ mục lá trong Merkle Tree tương ứng với tài khoản stealth cần bảo vệ.

    \item \textbf{Giai đoạn 2 -- Đề xuất Phục hồi (Guardian 1):}  Guardian 1 gọi \texttt{proposeRecovery(newRoot, newLeaf)}. Trước đó, server tính:
    \begin{align*}
        \texttt{oldLeaf} &= \texttt{leafMap}[i] \\
        \texttt{siblings} &= \texttt{createMerkleProof}(i) \\
        \texttt{newRoot} &= f_{\text{merkle}}(\texttt{newLeaf}, i, \texttt{siblings})
    \end{align*}
    và sinh bằng chứng ZK SMT Update qua endpoint \texttt{POST /leaves/recovery-proof}.

    \item \textbf{Giai đoạn 3 -- Bỏ phiếu Phê duyệt (Guardians 2..t):}  Các guardian còn lại gọi \texttt{approveRecovery(reqId)} để bỏ phiếu. Hợp đồng đảm bảo mỗi guardian chỉ bỏ phiếu một lần trên mỗi yêu cầu.

    \item \textbf{Giai đoạn 4 -- Thực thi Phục hồi:}  Khi \texttt{req.approvals} $\geq t$, gọi \texttt{executeRecovery(reqId, auth)}. Hợp đồng \texttt{SocialRecoveryModule} chuyển tiếp lệnh cập nhật cây với bằng chứng ZK:
    \[
        \texttt{treeManager.updateRoot(newRoot, newLeaf, mappedIndex, auth)}
    \]
    Hợp đồng \texttt{SMTUpdateVerifier} xác minh bằng chứng Groth16 trước khi cập nhật gốc on-chain.

    \item \textbf{Giai đoạn 5 -- Đồng bộ Server:}  Sau khi giao dịch on-chain được xác nhận, server đồng bộ cây cục bộ qua \texttt{POST /leaves/sync-recovery}.

    \item \textbf{Giai đoạn 6 -- Khôi phục Quyền chi tiêu:}  Chủ sở hữu mới dùng \texttt{spendPriv\_mới} để nhập ví và chi tiêu từ tài khoản stealth.
\end{enumerate}

\begin{figure}[H]
    \centering
    \includegraphics[width=0.8\textwidth]{Hinhve/usecase_social_recovery.png}
    \caption{Biểu đồ Use Case -- Chức năng Phục hồi xã hội (Social Recovery)}
    \label{fig:uc_social_recovery}
\end{figure}

\begin{figure}[H]
    \centering
    \includegraphics[width=1\textwidth]{Hinhve/class_social_recovery.png}
    \caption{Biểu đồ Lớp -- Chức năng Phục hồi xã hội}
    \label{fig:class_social_recovery}
\end{figure}

\begin{figure}[H]
    \centering
    \includegraphics[width=1\textwidth]{Hinhve/social_recovery_flow.png}
    \caption{Biểu đồ tuần tự UML -- Luồng phục hồi xã hội (Social Recovery)}
    \label{fig:sd_social_recovery}
\end{figure}

\subsection{Thiết kế Mạch Zero-Knowledge (ZK Circuits)}
\label{subsec:zk_circuits_design}

Hai mạch ZK-SNARK đóng vai trò cốt lõi trong việc đảm bảo tính bảo mật và riêng tư của hệ thống được phát triển bằng ngôn ngữ Circom:

\begin{itemize}
    \item \textbf{Mạch Chi tiêu (\texttt{stealth\_spend.circom}):} 
    Mạch này chứng minh người dùng sở hữu khóa bí mật tương ứng với một địa chỉ stealth hợp lệ nằm trong Merkle Tree, đồng thời tạo ra một \texttt{nullifier} duy nhất để chống chi tiêu kép (double-spending).
    \begin{itemize}
        \item \textit{Tín hiệu đầu vào công khai (Public Inputs):} \texttt{root} (gốc của Merkle Tree hiện tại), \texttt{nullifierHash} (giá trị băm dùng để hủy, chống chi tiêu kép).
        \item \textit{Tín hiệu đầu vào bí mật (Private Inputs):} \texttt{spendPriv} (khóa bí mật chi tiêu), \texttt{sharedSecretHash} (giá trị băm của shared secret), \texttt{senderIndexCommitment} (cam kết chỉ mục), và \texttt{siblings} (các nút anh em trên đường dẫn Merkle).
        \item \textit{Logic ràng buộc:} Mạch kiểm tra chữ ký khóa, tính toán lại lá từ các đầu vào bí mật, và xác minh đường dẫn Merkle từ lá đó đến \texttt{root}. Đồng thời, tính toán và đối chiếu \texttt{nullifierHash}.
    \end{itemize}

    \item \textbf{Mạch Cập nhật Cây (\texttt{smt\_update.circom}):}
    Mạch này cho phép Relayer chứng minh việc cập nhật một lá mới vào cây Merkle là hợp lệ mà không cần thực hiện toàn bộ phép tính băm tốn kém trên hợp đồng thông minh.
    \begin{itemize}
        \item \textit{Tín hiệu đầu vào công khai (Public Inputs):} \texttt{oldRoot} (gốc cũ), \texttt{newRoot} (gốc mới sau khi cập nhật), \texttt{leaf} (giá trị lá mới), \texttt{index} (vị trí chèn).
        \item \textit{Tín hiệu đầu vào bí mật (Private Inputs):} \texttt{siblings} (các nút anh em trên đường dẫn Merkle từ \texttt{index} lên gốc).
        \item \textit{Logic ràng buộc:} Mạch xác minh tính hợp lệ của đường dẫn Merkle từ vị trí \texttt{index} (đang rỗng hoặc giá trị cũ) đến \texttt{oldRoot}. Sau đó, thay thế giá trị tại \texttt{index} bằng \texttt{leaf} và tính toán lại đường dẫn để đảm bảo khớp với \texttt{newRoot}.
    \end{itemize}
\end{itemize}

% ═══════════════════════════════════════════════════════════════════════════════
\section{Xây dựng ứng dụng}

\subsection{Thư viện và công cụ sử dụng}
Quá trình xây dựng ứng dụng Ví Stealth Address được phát triển trên nhiều nền tảng và ngôn ngữ khác nhau, từ xử lý logic blockchain đến giao diện web và các thuật toán toán học Zero-Knowledge. Danh sách các thư viện và công cụ được liệt kê chi tiết trong Bảng \ref{table:tools}.

\begin{table}[H]
\centering{}
\resizebox{\textwidth}{!}{
    \begin{tabular}{|l|l|l|l|}
        \hline
        \textbf{Mục đích} & \textbf{Công cụ} & \textbf{Phiên bản} & \textbf{Địa chỉ URL} \\ \hline
        Ngôn ngữ lập trình Smart Contract & Solidity & v0.8.28 & https://soliditylang.org/ \\ \hline
        Khung làm việc (Framework) EVM & Hardhat & v2.22.x & https://hardhat.org/ \\ \hline
        Ngôn ngữ viết mạch ZK-SNARK & Circom & v2.0.0 & https://docs.circom.io/ \\ \hline
        Thư viện sinh và xác minh ZK Proof & SnarkJS & v0.7.0 & https://github.com/iden3/snarkjs \\ \hline
        Môi trường chạy Server Relayer & Node.js (Express) & v20.x & https://nodejs.org/ \\ \hline
        Thư viện giao tiếp Web3 / Mật mã & ethers.js & v6.11.1 & https://docs.ethers.org/v6/ \\ \hline
        Phát triển giao diện Frontend & React (Vite) & v18.2.0 & https://react.dev/ \\ \hline
    \end{tabular}
}
    \caption{Danh sách thư viện và công cụ sử dụng}
    \label{table:tools}
\end{table}

\subsection{Kết quả đạt được}
Kết quả thực nghiệm đã đóng gói thành công 4 phân hệ (module) chính của hệ thống, tương tác chặt chẽ với nhau theo mô hình Account Abstraction:
\begin{itemize}
    \item \textbf{Gói Blockchain (\texttt{blockchain/contracts}):} Đóng gói các hợp đồng thông minh lõi, bao gồm \texttt{StealthAccount.sol} (Ví AA), \texttt{ERC5564Announcer.sol} (Lưu vết giao dịch), \texttt{StealthTreeManager.sol} (Quản lý Merkle Tree on-chain) và \texttt{SocialRecoveryModule.sol} (Quản lý khôi phục xã hội).
    \item \textbf{Gói Frontend (\texttt{stealth-wallet}):} Giao diện người dùng Web Application, đảm nhiệm việc sinh khóa cục bộ (Local Key Generation), tự động quét các sự kiện trên chuỗi và tích hợp WebAssembly (WASM) để sinh bằng chứng ZK ngay trên trình duyệt mà không làm rò rỉ khóa riêng.
    \item \textbf{Gói Mạch ZK (\texttt{zk/circuits}):} Bao gồm các mạch \texttt{stealth.circom} và \texttt{smt\_update.circom} được biên dịch thành file WASM và zkey, phục vụ cho việc sinh và xác minh bằng chứng mật mã.
    \item \textbf{Gói Relayer Backend (\texttt{server}):} Server API viết bằng Node.js xử lý việc nhận các \texttt{UserOperation}, duy trì cơ sở dữ liệu cây Merkle off-chain (\texttt{leaves.json}) và thay mặt người dùng trả phí gas.
\end{itemize}

Thông tin thống kê về quy mô mã nguồn của ứng dụng được thể hiện qua Bảng \ref{table:stats}.

\begin{table}[H]
\centering{}
    \begin{tabular}{|l|c|c|c|}
        \hline
        \textbf{Phân hệ / Sản phẩm đóng gói} & \textbf{Số lượng tệp tin} & \textbf{Số dòng code (LOC)} & \textbf{Dung lượng mã nguồn} \\ \hline
        Smart Contracts (Solidity) & 12 & $\sim$ 1,500 & 85 KB \\ \hline
        Mạch ZK (Circom) & 5 & $\sim$ 500 & 20 KB \\ \hline
        Frontend Web App (React/JS) & 25 & $\sim$ 2,500 & 2.5 MB \\ \hline
        Relayer Server (Node.js) & 10 & $\sim$ 800 & 150 KB \\ \hline
        \textbf{Tổng cộng (Ước tính)} & \textbf{52} & \textbf{$\sim$ 5,300} & \textbf{$\sim$ 2.75 MB} \\ \hline
    \end{tabular}
    \caption{Thống kê quy mô mã nguồn ứng dụng}
    \label{table:stats}
\end{table}

\subsection{Minh họa các chức năng chính}

Phần này trình bày giao diện của các tính năng quan trọng nhất trong hệ thống kèm theo mô tả về luồng dữ liệu bên dưới.

\begin{figure}[H]
    \centering
    \includegraphics[width=0.8\textwidth]{Hinhve/placeholder_create_wallet.png}
    \caption{Giao diện khởi tạo Stealth Meta-Address}
    \label{fig:func1}
\end{figure}
\textbf{Giải thích:} Tại giao diện này, hệ thống sẽ sử dụng thuật toán đường cong elliptic \texttt{secp256k1} qua thư viện \texttt{ethers.js} để tạo hai cặp khóa: Khóa xem (Viewing Key) và Khóa chi tiêu (Spending Key). Hệ thống nối hai khóa công khai lại và mã hóa Base58 để hiển thị một \texttt{Stealth Meta-Address} duy nhất cho người dùng. Khóa cá nhân được mã hóa và lưu an toàn ở bộ nhớ cục bộ (Local Storage).

\begin{figure}[H]
    \centering
    \includegraphics[width=0.8\textwidth]{Hinhve/placeholder_send_transfer.png}
    \caption{Giao diện chuyển tiền ẩn danh (Stealth Transfer)}
    \label{fig:func2}
\end{figure}
\textbf{Giải thích:} Người gửi nhập Meta-Address của người nhận. Dưới nền tảng, ứng dụng tự động sinh một khóa tạm thời (Ephemeral Key), kết hợp với Khóa xem công khai của người nhận để sinh ra Bí mật chung (Shared Secret) theo giao thức ECDH. Địa chỉ ẩn danh cuối cùng được tạo ra và giao dịch được gửi lên chuỗi, đính kèm sự kiện \texttt{Announcement} ghi nhận Ephemeral Public Key.

\begin{figure}[H]
    \centering
    \includegraphics[width=0.8\textwidth]{Hinhve/placeholder_scanner.png}
    \caption{Giao diện quét và khám phá tài sản (Scanner)}
    \label{fig:func3}
\end{figure}
\textbf{Giải thích:} Giao diện cho phép người dùng khởi chạy tiến trình quét (sử dụng Web Worker để không gây giật lag trình duyệt). Trình quét tải toàn bộ sự kiện \texttt{Announcement} từ blockchain, kết hợp với Khóa xem bí mật để đối chiếu View Tag (1-byte). Nếu khớp, hệ thống giải mã để hiển thị số dư thực tế đang nằm trong các địa chỉ ẩn danh thuộc về người dùng.

\begin{figure}[H]
    \centering
    \includegraphics[width=0.8\textwidth]{Hinhve/placeholder_zk_spend.png}
    \caption{Giao diện rút tiền ẩn danh (ZK Spend)}
    \label{fig:func4}
\end{figure}
\textbf{Giải thích:} Khi người dùng yêu cầu rút tiền từ ví ẩn danh về ví sàn giao dịch, ứng dụng sẽ thu thập Merkle Root từ server, tính toán đường dẫn Merkle Path và chạy thư viện \texttt{snarkjs} (dưới dạng WASM) để sinh bằng chứng toán học. Toàn bộ quá trình tạo Proof chỉ tốn vài giây. Proof sau đó được gửi off-chain cho Relayer, Relayer sẽ trả phí gas và gửi giao dịch \texttt{execute} lên mạng lưới.

\section{Kiểm thử}
Quá trình kiểm thử tập trung vào hai nghiệp vụ lõi phức tạp nhất của hệ thống: Chi tiêu ZK (Stealth Spend) và Khôi phục mạng lưới xã hội (Social Recovery). Hệ thống sử dụng thư viện \texttt{Chai} và \texttt{Hardhat Testing} để viết unit test on-chain. Tổng cộng có 25 kịch bản kiểm thử (Test Cases) đã được chạy với tỷ lệ thành công (Pass rate) đạt 100\%.

Dưới đây là một số trường hợp kiểm thử (Test Cases) tiêu biểu.

\begin{table}[H]
\centering{}
\resizebox{\textwidth}{!}{
    \begin{tabular}{|p{1cm}|p{2cm}|p{4cm}|p{4cm}|p{3cm}|p{1.5cm}|}
        \hline
        \textbf{Mã TC} & \textbf{Chức năng} & \textbf{Các bước thực hiện} & \textbf{Kết quả mong đợi} & \textbf{Kết quả thực tế} & \textbf{Trạng thái} \\ \hline
        TC\_01 & ZK Spend & 1. Nhập đúng \texttt{spendPriv}. 2. Sinh ZK Proof. 3. Giao dịch \texttt{execute} tới Smart Contract. & \texttt{verifier.verifyProof} trả về \texttt{true}. Giao dịch chuyển ETH thành công. & Chuyển ETH thành công. & Pass \\ \hline
        TC\_02 & ZK Spend & 1. Nhập \texttt{spendPriv} sai (của người khác). 2. Sinh ZK Proof. 3. Giao dịch \texttt{execute}. & Proof không hợp lệ. Smart Contract revert lỗi \texttt{INVALID\_PROOF}. & Bị revert đúng lỗi chỉ định. & Pass \\ \hline
        TC\_03 & ZK Spend (Replay) & 1. Bắt gói tin ZK Proof cũ. 2. Đổi \texttt{indexCommitment} sang ví khác và gửi lại. & ZK Proof bị từ chối do không khớp tín hiệu Public Inputs gốc. & Bị revert lỗi tín hiệu không khớp. & Pass \\ \hline
        TC\_04 & Social Recovery & 1. 1 Guardian gọi \texttt{proposeRecovery}. 2. 2 Guardians gọi \texttt{approveRecovery} (Threshold=3). 3. Bất kỳ ai gọi \texttt{executeRecovery} kèm ZK Proof cập nhật Merkle. & Root mới được cập nhật trên \texttt{StealthTreeManager}. Ví thay đổi Owner thành công. & Cập nhật Root và Leaf thành công. & Pass \\ \hline
        TC\_05 & Social Recovery & 1. 1 Guardian gọi \texttt{proposeRecovery}. 2. 1 Guardian gọi \texttt{approveRecovery}. 3. Gọi \texttt{executeRecovery}. & Bị revert do số lượng phê duyệt (2) nhỏ hơn Threshold (3). & Bị revert lỗi \texttt{not enough approvals}. & Pass \\ \hline
    \end{tabular}
}
    \caption{Các trường hợp kiểm thử nghiệp vụ lõi (Test Cases)}
    \label{table:testcases}
\end{table}

\textbf{Tổng kết kiểm thử:} Quá trình kiểm tra đảm bảo hệ thống không tồn tại lỗ hổng rò rỉ khóa trong quá trình sinh bằng chứng, và bảo vệ tuyệt đối khỏi các dạng tấn công lặp lại (Replay Attacks) cũng như tấn công thay đổi trạng thái (State Manipulation) từ người bảo hộ giả mạo.

\section{Triển khai}
\subsection{Mô hình triển khai thử nghiệm}
Môi trường thực nghiệm được thiết lập chạy cục bộ (Localhost) để giả lập mạng lưới, nhằm phục vụ công tác tinh chỉnh và đo đạc thông số:
\begin{itemize}
    \item \textbf{Mạng Blockchain:} Triển khai trên node cục bộ \texttt{Hardhat Network} (ChainID: 31337). Quá trình deploy sinh ra 5 hợp đồng thông minh lõi và hệ thống tự động ghi địa chỉ hợp đồng vào biến môi trường của Frontend.
    \item \textbf{Server Backend:} Chạy trên nền tảng Node.js (cấu hình thử nghiệm giả lập: 4 Cores CPU, 8GB RAM). Server mở port 3001, lắng nghe yêu cầu khởi tạo cây Merkle và trả kết quả API dưới dạng JSON.
    \item \textbf{Trình duyệt Client:} Giao diện Web chạy qua Vite (port 5173). Thử nghiệm thực hiện trên trình duyệt Google Chrome v120.
\end{itemize}

\subsection{Kết quả và Thông số kỹ thuật đạt được}
Trong quá trình triển khai thử nghiệm, hệ thống đã ghi nhận các thông số hiệu năng (Performance Metrics) ở mức xuất sắc đối với một ứng dụng Web3 mật mã hóa cao:

\textbf{1. Hiệu năng tính toán ZK Proof ở Client:}
\begin{itemize}
    \item Độ sâu cây Merkle (Levels): 20 (Hỗ trợ hơn 1 triệu địa chỉ ẩn danh).
    \item Thời gian tải thư viện \texttt{snarkjs} và file WASM (2MB): $\sim 0.5$ giây.
    \item \textbf{Thời gian sinh bằng chứng (Generate Proof):} Khoảng $\sim 1.5 - 2.5$ giây. Đây là kết quả rất tốt, đảm bảo trải nghiệm người dùng (UX) không bị gián đoạn quá lâu.
\end{itemize}

\textbf{2. Chi phí Gas (Phí giao dịch On-chain):}
Mọi tính toán ZK được xác minh qua các Precompiled Contracts tích hợp sẵn trên EVM (như phép toán trên đường cong alt\_bn128/bn254).
\begin{itemize}
    \item Giao dịch \texttt{execute} (Xác minh ZK Spend + Chuyển tiền): $\sim 310,000$ Gas.
    \item Đề xuất khôi phục \texttt{proposeRecovery}: $\sim 95,000$ Gas.
    \item Thực thi khôi phục \texttt{executeRecovery} (Xác minh cập nhật SMT Merkle): $\sim 350,000$ Gas.
    \item \textit{Đánh giá:} Chi phí này rất cạnh tranh và tương đương với các giải pháp máy trộn (Mixer) nổi tiếng trên thị trường, nhưng mang lại mức độ phân quyền và bảo vệ định tuyến an toàn hơn.
\end{itemize}

\textbf{3. Tốc độ Quét (Scanning Speed):}
Nhờ kiến trúc \textit{View Tag Optimization} (chỉ đối chiếu byte đầu tiên của mã băm thay vì nhân ma trận toàn khối lượng giao dịch), tốc độ quét mạng lưới tăng gấp 255 lần.
\begin{itemize}
    \item Quét thử nghiệm 1,000 sự kiện \texttt{Announcement}: Xử lý cục bộ mất chưa đến 0.2 giây.
    \item Quét thử nghiệm 10,000 sự kiện: Chỉ mất $\sim 1.8$ giây, chứng minh hệ thống có khả năng chịu tải dữ liệu lịch sử cực lớn khi triển khai thực tế trên mạng Mainnet.
\end{itemize}

\end{document}